\section{金的地方统治：迁都、边防与北方生产}

\subsection{女真建国与东北资源}

账本将。若从，首先可见的不是一条已经完成的国家命令，而是道路、仓储、城镇、田土与人群被重新接入的过程。建国将地方联盟转入新政治框架 这一变化既有中枢决策的一面，也有地方必须筹粮、输送、登记、修治或避险的一面；因此不能把它写成只由宫廷人物推动的直线故事。

这一段的基本年序可由。它们的文体与观察位置并不相同：本纪或编年材料善于保存诏令、任免与军事节点，制度汇编更接近机构语言，碑刻、文书、遗址和地方材料则只在有限地点留下行动者的痕迹。把这些材料并排，不是为了拼出一幅无缝图景，而是为了分开日期、规范、实施与后果四种不同的判断。早金社会材料稀薄，不能由后出官制倒推

财政在这里不是抽象的‘国力’。一项边防、河工、迁都、互市或接收安排，至少牵动实物运输、货币或票据、胥吏文书与地方劳役中的若干环节。中央记录中出现的额数、户口、兵额或赋役名目，必须保留统计机关、地域和报告场景；它们不能直接换算为家庭的收入、损失或服役天数。相反，局部契约、墓志和考古发现也只能说明被保存下来的少数关系，不能反过来代表整个政权。

社会经验因而具有地域差异。城市市场与港口、边寨与河谷、寺院与家族、军户与流动者面对的并非同一种行政压力。材料中的官号、文字、籍属、宗教捐施或职业，也分别记录不同事实；不能用一个族名、一个制度名称或一座都城的经验将它们归并。把这种差异留在叙事中，才能解释同一政策为何会在不同地区产生不同的配合、拖延、规避或重新协商。

文化与宗教材料同样应放回制度环境。寺院题记、经卷、书院记录、学校条目或技术文本可以照亮书写、教育、捐施与知识传播的网络，却不自动证明国家控制了信仰，也不证明所有居民共享同一价值。它们与账本事件相接的方式，是展示政治、财政和交通如何创造某些行动机会，同时也展示现存证据无法覆盖的声音。

因此，阿骨打起兵并建金 的意义不在于提供一个可供道德裁断的孤立瞬间，而在于使、资源与社会关系暂时显形。后续研究仍须把同一事件放回相邻年度、对方政权材料和具体地理中核对；在缺少地方证据处，本书保留‘可见的是’与‘尚不能断定’之间的距离，而不以叙事流畅取代证据。

\subsection{燕京接收与华北道路}

账本将。若从，首先可见的不是一条已经完成的国家命令，而是道路、仓储、城镇、田土与人群被重新接入的过程。接收城市需要长期的人力和信息 这一变化既有中枢决策的一面，也有地方必须筹粮、输送、登记、修治或避险的一面；因此不能把它写成只由宫廷人物推动的直线故事。

这一段的基本年序可由。它们的文体与观察位置并不相同：本纪或编年材料善于保存诏令、任免与军事节点，制度汇编更接近机构语言，碑刻、文书、遗址和地方材料则只在有限地点留下行动者的痕迹。把这些材料并排，不是为了拼出一幅无缝图景，而是为了分开日期、规范、实施与后果四种不同的判断。城内交接与居民经验须由多类材料重建

财政在这里不是抽象的‘国力’。一项边防、河工、迁都、互市或接收安排，至少牵动实物运输、货币或票据、胥吏文书与地方劳役中的若干环节。中央记录中出现的额数、户口、兵额或赋役名目，必须保留统计机关、地域和报告场景；它们不能直接换算为家庭的收入、损失或服役天数。相反，局部契约、墓志和考古发现也只能说明被保存下来的少数关系，不能反过来代表整个政权。

社会经验因而具有地域差异。城市市场与港口、边寨与河谷、寺院与家族、军户与流动者面对的并非同一种行政压力。材料中的官号、文字、籍属、宗教捐施或职业，也分别记录不同事实；不能用一个族名、一个制度名称或一座都城的经验将它们归并。把这种差异留在叙事中，才能解释同一政策为何会在不同地区产生不同的配合、拖延、规避或重新协商。

文化与宗教材料同样应放回制度环境。寺院题记、经卷、书院记录、学校条目或技术文本可以照亮书写、教育、捐施与知识传播的网络，却不自动证明国家控制了信仰，也不证明所有居民共享同一价值。它们与账本事件相接的方式，是展示政治、财政和交通如何创造某些行动机会，同时也展示现存证据无法覆盖的声音。

因此，取燕京、灭辽并南侵宋境 的意义不在于提供一个可供道德裁断的孤立瞬间，而在于使、资源与社会关系暂时显形。后续研究仍须把同一事件放回相邻年度、对方政权材料和具体地理中核对；在缺少地方证据处，本书保留‘可见的是’与‘尚不能断定’之间的距离，而不以叙事流畅取代证据。

\subsection{刘豫政权与中原治理}

账本将。若从，首先可见的不是一条已经完成的国家命令，而是道路、仓储、城镇、田土与人群被重新接入的过程。代理和直辖各有征收与忠诚成本 这一变化既有中枢决策的一面，也有地方必须筹粮、输送、登记、修治或避险的一面；因此不能把它写成只由宫廷人物推动的直线故事。

这一段的基本年序可由。它们的文体与观察位置并不相同：本纪或编年材料善于保存诏令、任免与军事节点，制度汇编更接近机构语言，碑刻、文书、遗址和地方材料则只在有限地点留下行动者的痕迹。把这些材料并排，不是为了拼出一幅无缝图景，而是为了分开日期、规范、实施与后果四种不同的判断。政权更替不等于地方执行同时改变

财政在这里不是抽象的‘国力’。一项边防、河工、迁都、互市或接收安排，至少牵动实物运输、货币或票据、胥吏文书与地方劳役中的若干环节。中央记录中出现的额数、户口、兵额或赋役名目，必须保留统计机关、地域和报告场景；它们不能直接换算为家庭的收入、损失或服役天数。相反，局部契约、墓志和考古发现也只能说明被保存下来的少数关系，不能反过来代表整个政权。

社会经验因而具有地域差异。城市市场与港口、边寨与河谷、寺院与家族、军户与流动者面对的并非同一种行政压力。材料中的官号、文字、籍属、宗教捐施或职业，也分别记录不同事实；不能用一个族名、一个制度名称或一座都城的经验将它们归并。把这种差异留在叙事中，才能解释同一政策为何会在不同地区产生不同的配合、拖延、规避或重新协商。

文化与宗教材料同样应放回制度环境。寺院题记、经卷、书院记录、学校条目或技术文本可以照亮书写、教育、捐施与知识传播的网络，却不自动证明国家控制了信仰，也不证明所有居民共享同一价值。它们与账本事件相接的方式，是展示政治、财政和交通如何创造某些行动机会，同时也展示现存证据无法覆盖的声音。

因此，扶植刘豫后转为直接经营 的意义不在于提供一个可供道德裁断的孤立瞬间，而在于使、资源与社会关系暂时显形。后续研究仍须把同一事件放回相邻年度、对方政权材料和具体地理中核对；在缺少地方证据处，本书保留‘可见的是’与‘尚不能断定’之间的距离，而不以叙事流畅取代证据。

\subsection{中都与南侵的资源链}

账本将。若从，首先可见的不是一条已经完成的国家命令，而是道路、仓储、城镇、田土与人群被重新接入的过程。都城建设和远征共同挤压供给 这一变化既有中枢决策的一面，也有地方必须筹粮、输送、登记、修治或避险的一面；因此不能把它写成只由宫廷人物推动的直线故事。

这一段的基本年序可由。它们的文体与观察位置并不相同：本纪或编年材料善于保存诏令、任免与军事节点，制度汇编更接近机构语言，碑刻、文书、遗址和地方材料则只在有限地点留下行动者的痕迹。把这些材料并排，不是为了拼出一幅无缝图景，而是为了分开日期、规范、实施与后果四种不同的判断。战功、伤亡和军数均有战报夸张风险

财政在这里不是抽象的‘国力’。一项边防、河工、迁都、互市或接收安排，至少牵动实物运输、货币或票据、胥吏文书与地方劳役中的若干环节。中央记录中出现的额数、户口、兵额或赋役名目，必须保留统计机关、地域和报告场景；它们不能直接换算为家庭的收入、损失或服役天数。相反，局部契约、墓志和考古发现也只能说明被保存下来的少数关系，不能反过来代表整个政权。

社会经验因而具有地域差异。城市市场与港口、边寨与河谷、寺院与家族、军户与流动者面对的并非同一种行政压力。材料中的官号、文字、籍属、宗教捐施或职业，也分别记录不同事实；不能用一个族名、一个制度名称或一座都城的经验将它们归并。把这种差异留在叙事中，才能解释同一政策为何会在不同地区产生不同的配合、拖延、规避或重新协商。

文化与宗教材料同样应放回制度环境。寺院题记、经卷、书院记录、学校条目或技术文本可以照亮书写、教育、捐施与知识传播的网络，却不自动证明国家控制了信仰，也不证明所有居民共享同一价值。它们与账本事件相接的方式，是展示政治、财政和交通如何创造某些行动机会，同时也展示现存证据无法覆盖的声音。

因此，完颜亮政变、迁中都并准备南侵 的意义不在于提供一个可供道德裁断的孤立瞬间，而在于使、资源与社会关系暂时显形。后续研究仍须把同一事件放回相邻年度、对方政权材料和具体地理中核对；在缺少地方证据处，本书保留‘可见的是’与‘尚不能断定’之间的距离，而不以叙事流畅取代证据。

\subsection{世宗后的边区经济}

账本将。若从，首先可见的不是一条已经完成的国家命令，而是道路、仓储、城镇、田土与人群被重新接入的过程。和平时期仍有驻军、贸易和税役成本 这一变化既有中枢决策的一面，也有地方必须筹粮、输送、登记、修治或避险的一面；因此不能把它写成只由宫廷人物推动的直线故事。

这一段的基本年序可由。它们的文体与观察位置并不相同：本纪或编年材料善于保存诏令、任免与军事节点，制度汇编更接近机构语言，碑刻、文书、遗址和地方材料则只在有限地点留下行动者的痕迹。把这些材料并排，不是为了拼出一幅无缝图景，而是为了分开日期、规范、实施与后果四种不同的判断。一处边城或中都经验不能代替各路社会

财政在这里不是抽象的‘国力’。一项边防、河工、迁都、互市或接收安排，至少牵动实物运输、货币或票据、胥吏文书与地方劳役中的若干环节。中央记录中出现的额数、户口、兵额或赋役名目，必须保留统计机关、地域和报告场景；它们不能直接换算为家庭的收入、损失或服役天数。相反，局部契约、墓志和考古发现也只能说明被保存下来的少数关系，不能反过来代表整个政权。

社会经验因而具有地域差异。城市市场与港口、边寨与河谷、寺院与家族、军户与流动者面对的并非同一种行政压力。材料中的官号、文字、籍属、宗教捐施或职业，也分别记录不同事实；不能用一个族名、一个制度名称或一座都城的经验将它们归并。把这种差异留在叙事中，才能解释同一政策为何会在不同地区产生不同的配合、拖延、规避或重新协商。

文化与宗教材料同样应放回制度环境。寺院题记、经卷、书院记录、学校条目或技术文本可以照亮书写、教育、捐施与知识传播的网络，却不自动证明国家控制了信仰，也不证明所有居民共享同一价值。它们与账本事件相接的方式，是展示政治、财政和交通如何创造某些行动机会，同时也展示现存证据无法覆盖的声音。

因此，隆兴和议、世宗去世与嘉定和议 的意义不在于提供一个可供道德裁断的孤立瞬间，而在于使、资源与社会关系暂时显形。后续研究仍须把同一事件放回相邻年度、对方政权材料和具体地理中核对；在缺少地方证据处，本书保留‘可见的是’与‘尚不能断定’之间的距离，而不以叙事流畅取代证据。

\subsection{蒙古冲击下的地方收缩}

账本将。若从，首先可见的不是一条已经完成的国家命令，而是道路、仓储、城镇、田土与人群被重新接入的过程。防线收缩改变地方征发和迁徙 这一变化既有中枢决策的一面，也有地方必须筹粮、输送、登记、修治或避险的一面；因此不能把它写成只由宫廷人物推动的直线故事。

这一段的基本年序可由。它们的文体与观察位置并不相同：本纪或编年材料善于保存诏令、任免与军事节点，制度汇编更接近机构语言，碑刻、文书、遗址和地方材料则只在有限地点留下行动者的痕迹。把这些材料并排，不是为了拼出一幅无缝图景，而是为了分开日期、规范、实施与后果四种不同的判断。灭亡叙事不能掩盖不同地区的时间差

财政在这里不是抽象的‘国力’。一项边防、河工、迁都、互市或接收安排，至少牵动实物运输、货币或票据、胥吏文书与地方劳役中的若干环节。中央记录中出现的额数、户口、兵额或赋役名目，必须保留统计机关、地域和报告场景；它们不能直接换算为家庭的收入、损失或服役天数。相反，局部契约、墓志和考古发现也只能说明被保存下来的少数关系，不能反过来代表整个政权。

社会经验因而具有地域差异。城市市场与港口、边寨与河谷、寺院与家族、军户与流动者面对的并非同一种行政压力。材料中的官号、文字、籍属、宗教捐施或职业，也分别记录不同事实；不能用一个族名、一个制度名称或一座都城的经验将它们归并。把这种差异留在叙事中，才能解释同一政策为何会在不同地区产生不同的配合、拖延、规避或重新协商。

文化与宗教材料同样应放回制度环境。寺院题记、经卷、书院记录、学校条目或技术文本可以照亮书写、教育、捐施与知识传播的网络，却不自动证明国家控制了信仰，也不证明所有居民共享同一价值。它们与账本事件相接的方式，是展示政治、财政和交通如何创造某些行动机会，同时也展示现存证据无法覆盖的声音。

因此，蒙古攻金、失中都、汴京被围与蔡州陷落 的意义不在于提供一个可供道德裁断的孤立瞬间，而在于使、资源与社会关系暂时显形。后续研究仍须把同一事件放回相邻年度、对方政权材料和具体地理中核对；在缺少地方证据处，本书保留‘可见的是’与‘尚不能断定’之间的距离，而不以叙事流畅取代证据。

